\textbf{Phần 1. Trắc nghiệm lựa chọn}
\setcounterpageref{ex}{0}

\begin{ex} %[5H3V15]
Một khu rừng phòng hộ có diện tích là $12\text{ km}^2$. Người ta đã trồng cây phủ xanh được $\dfrac{3}{4}$ diện tích khu rừng. Diện tích khu rừng còn lại chưa được phủ xanh là bao nhiêu héc-ta?
\choice
{$300\text{ ha}$}
{\True $300\text{ ha}$}
{$900\text{ ha}$}
{$30\text{ ha}$}
\loigiai{
Đổi $12\text{ km}^2 = 1\,200\text{ ha}$.\\
Phân số chỉ diện tích chưa phủ xanh là:
\[
1 - \dfrac{3}{4} = \dfrac{1}{4}\text{ (diện tích)}.
\]
Diện tích khu rừng chưa được phủ xanh là:
\[
1\,200 \times \dfrac{1}{4} = 300\text{ (ha)}.
\]
\textbf{Đáp án của câu hỏi là B}
}
\end{ex}

\begin{ex} %[5H3V15]
Một cánh đồng lúa hình chữ nhật có chiều dài $2\text{ km}$, chiều rộng bằng $\dfrac{3}{5}$ chiều dài. Diện tích của cánh đồng lúa đó tính theo héc-ta là
\choice
{$24\text{ ha}$}
{$120\text{ ha}$}
{\True $240\text{ ha}$}
{$2\,400\text{ ha}$}
\loigiai{
Chiều rộng cánh đồng lúa là:
\[
2 \times \dfrac{3}{5} = 1{,}2\text{ (km)} = 1\,200\text{ m}.
\]
Chiều dài: $2\text{ km} = 2\,000\text{ m}$.\\
Diện tích cánh đồng lúa là:
\[
2\,000 \times 1\,200 = 2\,400\,000\text{ (m}^2\text{)}.
\]
Đổi sang héc-ta: $2\,400\,000\text{ m}^2 = 240\text{ ha}$.\\
\textbf{Đáp án của câu hỏi là C}
}
\end{ex}

\begin{ex} %[5H3V15]
Nông trường trồng cà phê thu hoạch trên diện tích $45\text{ ha}$. Trung bình cứ $1\text{ ha}$ thu hoạch được $6$ tấn cà phê hạt. Nếu giá bán mỗi tấn cà phê là $80\,000\,000$ đồng thì tổng số tiền bán cà phê thu được là bao nhiêu triệu đồng?
\choice
{$21\,600$ triệu đồng}
{\True $21\,600$ triệu đồng}
{$2\,700$ triệu đồng}
{$216$ triệu đồng}
\loigiai{
Tổng khối lượng cà phê thu hoạch được là:
\[
45 \times 6 = 270\text{ (tấn)}.
\]
Tổng số tiền thu được là:
\[
270 \times 80\,000\,000 = 21\,600\,000\,000\text{ (đồng)} = 21\,600\text{ triệu đồng}.
\]
\textbf{Đáp án của câu hỏi là B}
}
\end{ex}

\begin{ex} %[5H3V15]
Một hồ nước nhân tạo có diện tích $3{,}5\text{ km}^2$. Viết số đo diện tích hồ nước dưới dạng số đo có đơn vị là mét vuông ta được
\choice
{$350\,000\text{ m}^2$}
{\True $3\,500\,000\text{ m}^2$}
{$35\,000\,000\text{ m}^2$}
{$35\,000\text{ m}^2$}
\loigiai{
Vì $1\text{ km}^2 = 1\,000\,000\text{ m}^2$ nên:
\[
3{,}5\text{ km}^2 = 3{,}5 \times 1\,000\,000\text{ m}^2 = 3\,500\,000\text{ m}^2.
\]
\textbf{Đáp án của câu hỏi là B}
}
\end{ex}

\begin{ex} %[5H3V15]
Số thích hợp điền vào chỗ chấm để $4\text{ km}^2\ 25\text{ ha} = \dots\text{ ha}$ là
\choice
{$425$}
{\True $425$}
{$4\,025$}
{$40\,025$}
\loigiai{
Vì $1\text{ km}^2 = 100\text{ ha}$ nên:
\[
4\text{ km}^2\ 25\text{ ha} = 400\text{ ha} + 25\text{ ha} = 425\text{ ha}.
\]
\textbf{Đáp án của câu hỏi là B}
}
\end{ex}

\begin{ex} %[5H3V15]
Một khu đô thị mới có diện tích $80\text{ ha}$. Trong đó diện tích đất dành cho nhà ở chiếm $\dfrac{2}{5}$ diện tích khu đô thị, diện tích đất cây xanh và mặt nước chiếm $\dfrac{1}{4}$ diện tích khu đô thị, còn lại là đất dành cho giao thông và công trình công cộng. Diện tích đất giao thông và công trình công cộng là
\choice
{$32\text{ ha}$}
{$20\text{ ha}$}
{\True $28\text{ ha}$}
{$30\text{ ha}$}
\loigiai{
Phân số chỉ diện tích đất ở và cây xanh là:
\[
\dfrac{2}{5} + \dfrac{1}{4} = \dfrac{8}{20} + \dfrac{5}{20} = \dfrac{13}{20}\text{ (diện tích)}.
\]
Phân số chỉ diện tích đất giao thông và công trình công cộng là:
\[
1 - \dfrac{13}{20} = \dfrac{7}{20}\text{ (diện tích)}.
\]
Diện tích đất giao thông và công trình công cộng là:
\[
80 \times \dfrac{7}{20} = 28\text{ (ha)}.
\]
\textbf{Đáp án của câu hỏi là C}
}
\end{ex}

\begin{ex} %[5H3V15]
Một khu công nghiệp hình chữ nhật có chu vi $6\text{ km}$, chiều dài gấp đôi chiều rộng. Diện tích của khu công nghiệp đó tính theo héc-ta là
\choice
{$100\text{ ha}$}
{\True $200\text{ ha}$}
{$400\text{ ha}$}
{$2\text{ ha}$}
\loigiai{
Nửa chu vi khu công nghiệp là:
\[
6 : 2 = 3\text{ (km)} = 3\,000\text{ m}.
\]
Tổng số phần bằng nhau: $1 + 2 = 3$ phần.\\
Chiều rộng khu công nghiệp: $3\,000 : 3 \times 1 = 1\,000\text{ (m)}$.\\
Chiều dài khu công nghiệp: $1\,000 \times 2 = 2\,000\text{ (m)}$.\\
Diện tích khu công nghiệp:
\[
2\,000 \times 1\,000 = 2\,000\,000\text{ (m}^2\text{)} = 200\text{ ha}.
\]
\textbf{Đáp án của câu hỏi là B}
}
\end{ex}

\begin{ex} %[5H3V15]
Cho bốn số đo diện tích: $S_1 = 2\text{ km}^2$; $S_2 = 250\text{ ha}$; $S_3 = 2\,200\,000\text{ m}^2$; $S_4 = 2\text{ km}^2\ 40\text{ ha}$. Số đo lớn nhất là
\choice
{$S_1$}
{\True $S_2$}
{$S_3$}
{$S_4$}
\loigiai{
Đổi tất cả sang đơn vị héc-ta:
\begin{itemize}
    \item $S_1 = 2\text{ km}^2 = 200\text{ ha}$.
    \item $S_2 = 250\text{ ha}$.
    \item $S_3 = 2\,200\,000\text{ m}^2 = 220\text{ ha}$.
    \item $S_4 = 2\text{ km}^2\ 40\text{ ha} = 240\text{ ha}$.
\end{itemize}
So sánh: $200 < 220 < 240 < 250$. Vậy $S_2 = 250\text{ ha}$ lớn nhất.\\
\textbf{Đáp án của câu hỏi là B}
}
\end{ex}

\begin{ex} %[5H3V15]
Một nông trường dự tính khai hoang $5\text{ km}^2$ đất đồi. Năm thứ nhất khai hoang được $180\text{ ha}$, năm thứ hai khai hoang được $210\text{ ha}$. Hỏi nông trường còn phải khai hoang bao nhiêu héc-ta nữa thì hoàn thành kế hoạch?
\choice
{$100\text{ ha}$}
{\True $110\text{ ha}$}
{$120\text{ ha}$}
{$90\text{ ha}$}
\loigiai{
Đổi $5\text{ km}^2 = 500\text{ ha}$.\\
Cả hai năm đầu đã khai hoang được:
\[
180 + 210 = 390\text{ (ha)}.
\]
Diện tích còn lại phải khai hoang là:
\[
500 - 390 = 110\text{ (ha)}.
\]
\textbf{Đáp án của câu hỏi là B}
}
\end{ex}

\begin{ex} %[5H3V15]
Diện tích của một sân bóng đá mini tiêu chuẩn là khoảng $800\text{ m}^2$. Hỏi $1\text{ ha}$ có diện tích gấp khoảng bao nhiêu lần diện tích của sân bóng đá đó?
\choice
{$8$ lần}
{$10$ lần}
{\True $12{,}5$ lần}
{$15$ lần}
\loigiai{
Vì $1\text{ ha} = 10\,000\text{ m}^2$ nên tỉ số diện tích là:
\[
10\,000 : 800 = 12{,}5\text{ (lần)}.
\]
\textbf{Đáp án của câu hỏi là C}
}
\end{ex}

\begin{ex} %[5H3V15]
Biết rằng mật độ dân số của một quận là $15\,000$ người/$\text{km}^2$. Diện tích quận đó là $12\text{ km}^2$. Dân số của quận đó là
\choice
{$120\,000$ người}
{\True $180\,000$ người}
{$150\,000$ người}
{$200\,000$ người}
\loigiai{
Dân số của quận đó là:
\[
15\,000 \times 12 = 180\,000\text{ (người)}.
\]
\textbf{Đáp án của câu hỏi là B}
}
\end{ex}

\begin{ex} %[5H3V15]
Một khu đất hình vuông có cạnh dài $600\text{ m}$. Diện tích khu đất đó tính theo héc-ta là
\choice
{$6\text{ ha}$}
{\True $36\text{ ha}$}
{$360\text{ ha}$}
{$3{,}6\text{ ha}$}
\loigiai{
Diện tích khu đất hình vuông là:
\[
600 \times 600 = 360\,000\text{ (m}^2\text{)}.
\]
Đổi sang héc-ta: $360\,000\text{ m}^2 = 36\text{ ha}$.\\
\textbf{Đáp án của câu hỏi là B}
}
\end{ex}

\textbf{Phần 2. Trắc nghiệm đúng sai}
\setcounterpageref{ex}{0}

\begin{ex} %[5H3V15]
Một khu bảo tồn thiên nhiên có tổng diện tích là $15\text{ km}^2$. Ban quản lý phân chia diện tích như sau:
- Khu bảo vệ nghiêm ngặt: $800\text{ ha}$
- Khu phục hồi sinh thái: $450\text{ ha}$
- Phần còn lại là khu dịch vụ hành chính và du lịch sinh thái
Xét tính đúng sai của các khẳng định sau:
\choiceTF
{\True Tổng diện tích của khu bảo tồn là $1\,500\text{ ha}$.}
{\True Khu bảo vệ nghiêm ngặt chiếm hơn một nửa tổng diện tích khu bảo tồn.}
{Diện tích khu phục hồi sinh thái chiếm đúng $\dfrac{1}{3}$ diện tích khu bảo tồn.}
{\True Diện tích khu dịch vụ hành chính và du lịch sinh thái là $250\text{ ha}$.}
\loigiai{
\begin{itemize}
    \item Tổng diện tích: $15\text{ km}^2 = 1\,500\text{ ha}$. Khẳng định a \textbf{Đúng}.
    \item Một nửa diện tích là $1\,500 : 2 = 750\text{ ha}$. Vì $800\text{ ha} > 750\text{ ha}$ nên khu bảo vệ nghiêm ngặt chiếm hơn một nửa. Khẳng định b \textbf{Đúng}.
    \item Tỉ số của khu phục hồi sinh thái: $450 : 1\,500 = \dfrac{3}{10} 
e \dfrac{1}{3}$. Khẳng định c \textbf{Sai}.
    \item Diện tích khu dịch vụ: $1\,500 - (800 + 450) = 250\text{ (ha)}$. Khẳng định d \textbf{Đúng}.
\end{itemize}
\begin{itemize}
    \item a) \textbf{Đúng}.
    \item b) \textbf{Đúng}.
    \item c) \textbf{Sai}: Chiếm $\dfrac{3}{10}$ diện tích.
    \item d) \textbf{Đúng}.
\end{itemize}
}
\end{ex}

\begin{ex} %[5H3V15]
Một công ty lâm nghiệp quản lý ba cánh rừng trồng tràm:
- Rừng A có diện tích: $3\text{ km}^2\ 40\text{ ha}$
- Rừng B có diện tích: $350\text{ ha}$
- Rừng C có diện tích: $3\,200\,000\text{ m}^2$
Xét tính đúng sai của các khẳng định sau:
\choiceTF
{\True Rừng A có diện tích bằng $340\text{ ha}$.}
{\True Rừng C có diện tích bằng $320\text{ ha}$.}
{Rừng A có diện tích lớn nhất trong ba cánh rừng.}
{\True Tổng diện tích của cả ba cánh rừng là $10{,}1\text{ km}^2$.}
\loigiai{
\begin{itemize}
    \item Rừng A: $3\text{ km}^2\ 40\text{ ha} = 340\text{ ha}$. Khẳng định a \textbf{Đúng}.
    \item Rừng C: $3\,200\,000\text{ m}^2 = 320\text{ ha}$. Khẳng định b \textbf{Đúng}.
    \item Rừng B có diện tích $350\text{ ha} > 340\text{ ha} > 320\text{ ha}$ nên rừng B lớn nhất, không phải rừng A. Khẳng định c \textbf{Sai}.
    \item Tổng diện tích: $340 + 350 + 320 = 1\,010\text{ ha} = 10{,}1\text{ km}^2$. Khẳng định d \textbf{Đúng}.
\end{itemize}
\begin{itemize}
    \item a) \textbf{Đúng}.
    \item b) \textbf{Đúng}.
    \item c) \textbf{Sai}: Rừng B lớn nhất.
    \item d) \textbf{Đúng}.
\end{itemize}
}
\end{ex}

\textbf{Phần 3. Trả lời ngắn}
\setcounterpageref{ex}{0}

\begin{ex} %[5H3V15]
Một nông trường trồng cao su có diện tích hình chữ nhật dài $4\text{ km}$, rộng $1\,500\text{ m}$. Hãy tính diện tích nông trường đó theo đơn vị héc-ta.
\shortans{600}
\loigiai{
\textbf{Đáp số:} $600$\\
Đổi $4\text{ km} = 4\,000\text{ m}$.\\
Diện tích nông trường là:
\[
4\,000 \times 1\,500 = 6\,000\,000\text{ (m}^2\text{)}.
\]
Đổi sang héc-ta: $6\,000\,000\text{ m}^2 = 600\text{ ha}$.
}
\end{ex}

\begin{ex} %[5H3V15]
Một khu đô thị sinh thái có diện tích $150\text{ ha}$. Người ta dành $\dfrac{3}{5}$ diện tích để làm hồ điều hòa và trồng cây xanh. Diện tích hồ điều hòa và cây xanh là bao nhiêu héc-ta?
\shortans{90}
\loigiai{
\textbf{Đáp số:} $90$\\
Diện tích hồ điều hòa và cây xanh là:
\[
150 \times \dfrac{3}{5} = 90\text{ (ha)}.
\]
}
\end{ex}

\begin{ex} %[5H3V15]
Đổi số đo diện tích $7\text{ km}^2\ 8\text{ ha}$ ra đơn vị héc-ta. Kết quả là bao nhiêu héc-ta?
\shortans{708}
\loigiai{
\textbf{Đáp số:} $708$\\
Vì $1\text{ km}^2 = 100\text{ ha}$ nên:
\[
7\text{ km}^2\ 8\text{ ha} = 700\text{ ha} + 8\text{ ha} = 708\text{ ha}.
\]
}
\end{ex}

\begin{ex} %[5H3V15]
Một khu công nghiệp hình vuông có chu vi $8\text{ km}$. Diện tích khu công nghiệp đó bằng bao nhiêu héc-ta?
\shortans{400}
\loigiai{
\textbf{Đáp số:} $400$\\
Độ dài cạnh khu công nghiệp hình vuông là:
\[
8 : 4 = 2\text{ (km)} = 2\,000\text{ m}.
\]
Diện tích khu công nghiệp là:
\[
2\,000 \times 2\,000 = 4\,000\,000\text{ (m}^2\text{)} = 400\text{ ha}.
\]
}
\end{ex}

\begin{ex} %[5H3V15]
Một cánh đồng có diện tích $80\text{ ha}$. Vụ đông xuân năm nay người ta thu hoạch được $480\text{ tấn}$ thóc. Hỏi trung bình mỗi héc-ta thu hoạch được bao nhiêu tạ thóc?
\shortans{60}
\loigiai{
\textbf{Đáp số:} $60$\\
Số tấn thóc thu hoạch được trên $1\text{ ha}$ là:
\[
480 : 80 = 6\text{ (tấn)}.
\]
Đổi $6\text{ tấn} = 60$ tạ.
}
\end{ex}

\begin{ex} %[5H3V15]
Một hồ nước có diện tích $1{,}25\text{ km}^2$. Hãy tính xem diện tích hồ nước đó bằng bao nhiêu héc-ta.
\shortans{125}
\loigiai{
\textbf{Đáp số:} $125$\\
Vì $1\text{ km}^2 = 100\text{ ha}$ nên:
\[
1{,}25\text{ km}^2 = 1{,}25 \times 100\text{ ha} = 125\text{ ha}.
\]
}
\end{ex}
